
\subsection{Progressive Viterbi reconstruction (ProgRec)}

This algorithm is performed at each sibling pair to reconstruct the parent, in the usual manner of progressive alignment.
It begins with the construction of the Triad Pair HMM giving the distribution of a node and its two children.

The Pair HMM construction algorithm takes as input:
\begin{enumerate}
\item For each of two branches $e \in \{ 1, 2 \}$ leading to node $v$:
  an \emph{order-1 WFST} $W_e$ specified by its transition weight tensor
  \begin{equation}
  w_e(s, \anctok, \destok, s', \anctok', \destok')
  \label{eq:wfst-tensor}
  \end{equation}
  where $s, s' \in \smide$ are the 5-state pair HMM states,
  $\anctok, \anctok' \in \alphabet \cup \{\epsilon\}$ are ancestor characters
  (with $\epsilon$ denoting BOS/no-context),
  and $\destok, \destok' \in \alphabet \cup \{\epsilon\}$ are descendant characters.
  The tensor encodes the weight of transitioning from state $(s, \anctok, \destok)$
  to $(s', \anctok', \destok')$.

\item An \emph{order-1 singlet HMM} $R$ above the root, specified by
  \begin{equation}
  r(s, \anctok, s', \anctok')
  \label{eq:singlet-tensor}
  \end{equation}
  where $s, s' \in \{\sta, \ins, \fin\}$ and $\anctok, \anctok' \in \alphabet \cup \{\epsilon\}$.
\end{enumerate}

The WFSTs and singlet HMM are decoupled from their construction method
(algebraic distillation, numeric distillation, direct TKF92 parameterization, etc.).
They may have been produced by the distillation procedures of
Sections~\ref{sec:distill-hmm}--\ref{sec:order1trans},
by parametric construction using a simpler model (e.g. TKF92),
or by any other method that yields valid weight tensors.

\paragraph{State semantics.}
Each branch WFST state is a triple $(s, \anctok, \destok)$ where:
\begin{itemize}
\item $\mat$: consumes one ancestor character and emits one descendant character.
  Updates both $\anctok \to \anctok'$ and $\destok \to \destok'$.
\item $\ins$: emits one descendant character without consuming ancestor input.
  Updates $\destok \to \destok'$; preserves $\anctok' = \anctok$.
\item $\del$: consumes one ancestor character without emitting descendant output.
  Updates $\anctok \to \anctok'$; preserves $\destok' = \destok$.
\item $\sta$: start state (both contexts are $\epsilon$).
\item $\fin$: end state.
\end{itemize}
The singlet HMM emits only on transitions into the $\ins$ states. As with the branch WFSTs, the emitted character is uniquely determined by the destination state.

Let $I_e(s, \anctok, \destok, s', \anctok', \destok')$ be the indicator for whether the WFST transition
$(s, \anctok, \destok) \to (s', \anctok', \destok')$ is allowed using the above state semantics.
Let $I_s(s, \anctok, s', \anctok')$ be the corresponding indicator for whether the singlet transition
$(s, \anctok) \to (s', \anctok')$ is allowed.

\paragraph{Triad Pair HMM.}
The Triad Pair HMM represents the composition of the singlet HMM with the intersection of the two WFSTs.
It has states of the form $(s, \xstate_1, \xstate_2, \anctok, \destok_1, \destok_2)$
where $s$ represents the singlet state,
$\xstate_e$ represents the state of branch $e$'s WFST,
$\anctok$ represents the ancestral context (the last-emitted ancestral character, or $\epsilon$ for BOS/no-context),
$\destok_e$ represents the $e$-descendant context.

Only some states are accessible:
\begin{itemize}
\item The allowed $(s, \xstate_1, \xstate_2)$ state combinations are $\{\sta\sta\sta, \sta\sta\ins, \sta\ins\sta, \sta\ins\ins, \ins\mat\mat, \ins\mat\ins, \ins\mat\del, \ins\del\mat, \ins\del\ins, \ins\del\del, \ins\ins\mat, \ins\ins\ins, \ins\ins\del, \fin\fin\fin\}$.
\item The ancestral context cannot be empty ($\epsilon$) unless the branch state is one of $\sta\sta\sta$, $\sta\sta\ins$, $\sta\ins\sta$, $\sta\ins\ins$, or $\fin\fin\fin$.
\item The 1-descendant context cannot be empty unless the branch state is one of $\sta\sta\sta$, $\sta\sta\ins$, $\ins\del\mat$, $\ins\del\del$, $\ins\del\ins$, or $\fin\fin\fin$.
\item The 2-descendant context cannot be empty unless the branch state is one of $\sta\sta\sta$, $\sta\ins\sta$, $\ins\mat\del$, $\ins\del\del$, $\ins\ins\del$, or $\fin\fin\fin$.
\end{itemize}

\paragraph{State update ordering.}
We impose an ordering on state updates to avoid double-counting insert paths.
Let $\delta_e(s, \xstate_1, \xstate_2)$ be the indicator for whether a branch transducer changes state on entry into state $(s, \xstate_1, \xstate_2 \ldots)$.
Let $\delta_0(s, \xstate_1, \xstate_2)$ be the indicator for whether the root singlet HMM changes state.
\begin{eqnarray*}
    \delta_0(s, \xstate_1, \xstate_2) & = & \delta(\xstate_1 \neq \ins) \delta(\xstate_2 \neq \ins) \\
    \delta_1(s, \xstate_1, \xstate_2) & = & \delta(\xstate_2 \neq \ins) + \delta(\xstate_1 = \ins) \delta(\xstate_2 = \ins) \\
    \delta_2(s, \xstate_1, \xstate_2) & = & \delta(\xstate_1 \neq \ins)
\end{eqnarray*}

The transition matrix of the Triad Pair HMM has the following form
\[
w_{12}(s, \xstate_1, \xstate_2, \anctok, \destok_1, \destok_2, s', \xstate_1', \xstate_2', \anctok', \destok_1', \destok_2')
=
\left( w_0(s, \anctok, s', \anctok') I_0(s, \anctok, s', \anctok') \right)^{\delta_0(s', \xstate_1', \xstate_2')}
\delta(s = s', \anctok = \anctok')^{1 - \delta_0(s', \xstate_1', \xstate_2')}
\prod_{e \in \{1,2\}}
\left( w_e(\xstate_e, \anctok, \destok_e, \xstate_e', \anctok', \anctok', \destok_e') I_e(\xstate_e, \anctok, \destok_e, \xstate_e', \anctok', \anctok', \destok_e') \right)^{\delta_e(s', \xstate_1', \xstate_2')}
\delta(\xstate_e = \xstate_e', \destok_e = \destok_e')^{1 - \delta_e(s', \xstate_1', \xstate_2')}
\]

\paragraph{Progressive reconstruction.}
The Triad Pair HMM then can be used for Viterbi alignment of two siblings. The maximum likelihood parent sequence can be read off from the Viterbi state path.

\paragraph{Profile-based reconstruction.}
For the special case of models such as TKF91 and TKF92 where the indel and substitution processes are decoupled,
the order-1 structure of the models reduces to a simpler product of context-independent transitions with state-associated emissions.
In such cases we can maintain the ancestral profiles as position-specific weight matrices (PSWMs) instead of explicit sequences.
